% Part: sets-functions-relations
% Chapter: functions
% Section: basics

\documentclass[../../../include/open-logic-section]{subfiles}

\begin{document}

\olfileid{sfr}{fun}{bas}
\olsection{De basis}

\begin{explain}
Een \emph{functie} koppelt elk !!{element} uit een bepaalde verzameling aan
één bepaald !!{element} uit een andere verzameling, of soms uit dezelfde.
Optellen met~$1$ geeft bijvoorbeeld een functie: bij elk getal~$n$ hoort
precies één getal~$n+1$.
  
De invoer van een functie kan ook een paar, een drietal of iets met nog meer
delen zijn. Daar geeft de functie dan een uitvoer bij. Veel functies kennen we
al van gewoon rekenen.
Optellen en vermenigvuldigen zijn er voorbeelden van: je voert twee getallen
in en krijgt één getal als uitkomst.

Als wiskundig, abstract begrip is een functie een \emph{zwarte doos}. Het gaat
er alleen om welke uitvoer bij welke invoer hoort. Hoe je die uitvoer berekent,
doet er niet toe.
\end{explain}

\begin{defn}[Functie]
Een \emph{functie} $f \colon A \to B$ koppelt elk !!{element} van~$A$ aan
precies één !!{element} van~$B$.

$A$ heet het \emph{domein} van~$f$: daar komen de invoeren vandaan. $B$ heet
het \emph{codomein} van~$f$: daarin liggen de mogelijke uitkomsten. De
!!{element}en van~$A$ heten invoeren of \emph{argumenten}
van~$f$. Het !!{element} van~$B$ dat bij een argument~$x$ hoort volgens~$f$,
heet de \emph{waarde van~$f$} bij argument~$x$. We schrijven die als~$f(x)$.

Het \emph{bereik} $\ran{f}$ van~$f$ is het deel van het codomein waarin de
waarden van~$f$ staan die bij minstens één argument horen; $\ran{f} =
\Setabs{f(x)}{x \in A}$.
\end{defn}

Het plaatje in \olref{fig:function} kan helpen. De ellips links is het
\emph{domein} van de functie en de ellips rechts is het \emph{codomein}.
Een pijl gaat van een \emph{argument} in het domein naar de \emph{waarde} die
er in het codomein bij hoort.

\begin{figure}
  \olasset{assets/diagrams/function.tikz}
  \caption{Een functie koppelt elk !!{element} van de ene verzameling aan
    !!a{element} van een andere. Een pijl gaat van een argument in het
    domein naar de waarde die er in het codomein bij hoort.}
  \ollabel{fig:function}
\end{figure}

\begin{ex}
Bij vermenigvuldigen voer je twee natuurlijke getallen in en krijg je een
natuurlijk getal terug. Dit is dus een functie van $\Nat \times \Nat$ (het
domein) naar $\Nat$ (het codomein). Het bereik is ook $\Nat$, want elke
$n \in \Nat$ is gelijk aan $n \times 1$.
\end{ex}

\begin{ex}
Vermenigvuldigen is een functie: bij elke invoer---elk paar natuurlijke
getallen---hoort maar één uitvoer: $\times \colon \Nat^2 \to
\Nat$. Koppel je een natuurlijk getal~$n$ aan een reëel getal~$x$ waarvoor
$x^2 = n$, dan krijg je geen functie. Elk positief geheel getal~$n$ heeft
namelijk twee vierkantswortels: $\sqrt{n}$ en~$-\sqrt{n}$. We kunnen er een
functie van maken door alleen de niet-negatieve vierkantswortel terug te geven:
$\sqrt{\phantom{X}} \colon \Nat \to \Real$.
\end{ex}

\begin{ex}
Koppel je elke leerling in een klas aan het eindcijfer dat die leerling in die
klas krijgt, dan krijg je een functie. Een leerling kan binnen dezelfde klas
geen twee verschillende eindcijfers krijgen. Koppel je elke leerling aan de
ouders van die leerling, dan krijg je geen functie: een leerling kan nul, twee
of meer ouders hebben.
\end{ex}

\begin{explain}
Om precies te zeggen welke functie we bedoelen, moeten we voor elke mogelijke
invoer aangeven welke waarde erbij hoort. Dat kan met een formule, met een
manier om de waarde uit te rekenen of met een lijst waarin voor elke invoer de
waarde staat.
Welke manier we ook kiezen, bij elke invoer moet één, en maar één, waarde staan.
\end{explain}


\begin{ex}
Definieer $f \colon \Nat \to \Nat$ met $f(x) = x+1$. Daarmee hebben we een
functie $f$ die natuurlijke getallen binnenkrijgt en natuurlijke getallen
teruggeeft. Voer je een natuurlijk getal~$x$ in, dan geeft $f$ de opvolger,
het volgende getal~$x+1$. Het codomein $\Nat$ is hier niet het bereik van~$f$:
het natuurlijke getal~$0$ komt nooit uit deze functie, want het volgt op geen enkel
natuurlijk getal. Het bereik van~$f$ bestaat uit alle positieve gehele
getallen, $\Int^{+}$.
\end{ex}

\begin{ex}\ollabel{examplefunext}
Definieer $g \colon \Nat \to \Nat$ met $g(x) = x+2-1$. Dit is een functie
$g$ die natuurlijke getallen binnenkrijgt en natuurlijke getallen teruggeeft.
Bij een natuurlijk getal~$x$ geeft $g$ de uitkomst die je krijgt door vanaf
$x$ twee stappen vooruit en één stap terug te gaan, dus $x+1$.
\end{ex}

\begin{explain}
We hebben nu twee functies, $f$ en $g$, met verschillende \emph{definities}.
Toch zijn het \emph{dezelfde functie}. Voor elk natuurlijk getal~$n$ geldt
namelijk $f(n) = n+1 = n+2-1 = g(n)$. Anders gezegd: de definities van $f$
en~$g$ gebruiken verschillende vergelijkingen, maar leggen dezelfde koppeling
vast. Daarachter zit het extensionaliteitsbeginsel voor functies: twee functies
zijn gelijk als ze bij elke invoer dezelfde waarde geven. In symbolen:
\[
  \text{als }\forall x\, f(x) = g(x)\text{, dan }f = g
\]
Dit geldt alleen als $f$ en~$g$ hetzelfde domein én hetzelfde codomein hebben.
\end{explain}

\begin{ex}
We kunnen een functie ook geval voor geval beschrijven: per geval geven we aan
welke regel geldt. Zo kunnen we $h \colon \Nat \to \Nat$ definiëren:
\[
h(x) =
\begin{cases}
  \frac{x}{2} & \text{als $x$ even is} \\
  \frac{x+1}{2} & \text{als $x$ oneven is.}
\end{cases}
\]
Elk natuurlijk getal is even of oneven, dus de uitvoer van deze functie is
altijd een natuurlijk getal. Let er bij zo'n beschrijving op dat elke mogelijke
invoer in precies één geval valt. Soms moet je bewijzen dat elke mogelijke
invoer in een van de gevallen valt en nooit in twee gevallen tegelijk.
\end{ex}
\end{document}
